\section{方法与协议}
\label{sec:method}

读数分四步。前三步都在同一张切片上完成，第四步换到单细胞样本均值，不把两套坐标混在一起。

\begin{figure}[htbp]
  \centering
  \includegraphics[width=0.92\textwidth]{fig_flow.pdf}
  \caption{四步协议把接触检验、界面信号和腹膜方向分开。右侧为每步的数据或判定，左侧为操作。}
  \label{fig:flow}
\end{figure}

左侧四格按顺序完成质控、150\,$\mu$m接触、界面中位数和腹膜样本均值（图~\ref{fig:flow}）。接触的对照是保持两类数目的标签打乱，界面信号的对照是坐标打乱，两侧都是200次。保留线写在第三步右侧：至少5/9例，且留一后仍至少5/8。第四步的3份腹膜肿瘤和26份原发肿瘤只比较均值，GC6-PM不进入9例排序。半径、分位和这对保留线在出结果之前已经写死，后面的表按这套线报告，不另换阈值。

邻域用点位的显微坐标，半径150\,$\mu$m，由Visium点直径55\,$\mu$m除以全分辨率图像中的点直径像素数换算。配体受体来自CellChat人类库里配体和受体都是单个基因符号的组合，两基因需在全部切片中出现，共1000对。另有5对事先列入：CCL2--CCR2、C3--C3AR1、SPP1--CD44、TGFB1--TGFBR1、TGFB1--TGFBR2。名称里含下划线的复合物没有评分。

在这个可加读数里，去掉一条边之后界面中位数下降的量，就是这条边自己的中位数。报告的检验是该中位数是否高于坐标打乱，而不是边与边争夺同一份通量之后的再分配。单细胞矩阵若呈整数计数，则按细胞做每1万的$\log(1+x)$再对细胞取均值；样本之间比较的是这个均值。

另一次计算不替换上面的保留线。每个点位用四个标志基因程序做非负最小二乘，得到成纤维、巨噬、上皮和T细胞比例，四项加总为1。界面权重是成纤维比例乘巨噬比例。切片读数是该权重下的信号加权平均。空分布把配体表达图平移200、300、400或500\,$\mu$m，方向随机，200次，种子20261004；平移后找不到60\,$\mu$m内真实点位的位置记为缺失。主半径仍是150\,$\mu$m，另报100、200、300\,$\mu$m。9例汇总是中位数，区间来自对9例做2000次有放回抽样。只计算事先点名的配体受体，加上TGFB1乘以TGFBR1与TGFBR2的较小值。不重新筛选1000对，也不再做5/9投票。

同一套四个程序用在GSE183904的细胞上。细胞归入均值最大、且该均值大于0并高于其余三类的那一类；否则不归类。一类少于20个细胞的样本，该类均值不进入中位数。腹膜肿瘤3份对原发肿瘤26份只比方向。

乘积是事后指定的。每份样本先取配体所在类的均值、受体所在类的均值，再相乘。统计量是3份腹膜中位数减去原发中位数。在该条乘积有限的样本里，穷举所有把3份标成腹膜的组合。单侧p是统计量大于等于观察值的组合比例。7条一起做Benjamini--Hochberg。sample38相对sample36只列差值。这个乘积不是CellChat的通讯概率，也不是把一条边从网络里删掉之后的下游变化。
